@ID: ON19D1P1U
@BIDANG: Analisis Real
Diberikan fungsi $f:\mathbb{R}\rightarrow\mathbb{R}$ memenuhi sifat bahwa jika $\forall\lambda\in[0,1]$ dan $x,y\in\mathbb{R}$ maka $\lambda f(x)+(1-\lambda)f(y)-f(\lambda x+(1-\lambda)y)\ge0$. Buktikan:
(a) Untuk semua $x$ dan $y$ dengan $x<y$ maka $f(\frac{x+y}{2})\le\frac{f(x)+f(y)}{2}$.
(b) $\int_{0}^{2\pi}f(x)\cos x\,dx$ adalah positif.

@ID: ON19D1P2U
@BIDANG: Struktur Aljabar
Misalkan $G$ adalah grup dan $H$ subgrup dari $G$. Buktikan bahwa tidak ada subgrup $S$ dari $G$ dengan $S\ne G$ dan $S\supseteq G-H$ .

@ID: ON19D1P3U
@BIDANG: Kombinatorika
Diberikan bilangan bulat $n$ dan $k$ dengan $0\le k\le n/2$. Tuliskan sebuah bukti kombinatorial dari persamaan

$$
\sum_{m=k}^{n-k}\binom{m}{m-k}\binom{n-m}{k}=\binom{n+1}{2k+1}
$$



@ID: ON19D1P4U
@BIDANG: Analisis Kompleks
Diberikan fungsi kompleks $f(z)=\sum_{n=0}^{\infty}c_{n}z^{n}$ yang bersifat analitik pada $\mathcal{U}:=\{z\in\mathbb{C}:|z|<1\}$ dan diasumsikan bahwa $W:=\iint_{\mathcal{U}}|f^{\prime}(z)|^{2}dxdy<\infty$. Buktikan:
(a) $W=\pi\sum_{n=1}^{\infty}n|c_{n}|^{2}$.
(b) untuk setiap $z\in\mathcal{U}$ berlaku $|f(z)-f(0)|\le\sqrt{\frac{W}{\pi}\ln(\frac{1}{1-|z|^{2}})}$.

@ID: ON19D1P5U
@BIDANG: Aljabar Linear
Diberikan $F$ dan $H$ matriks berukuran $n\times n$ dengan $HF-FH=2^{2019}F$. Misalkan $v$ adalah vektor eigen dari $H$ dengan $Fv\ne0$. Buktikan terdapat bilangan bulat positif $k$ sehingga $F^{k}v$ adalah vektor eigen dari $F$.

@ID: ON19D2P1U
@BIDANG: Kombinatorika
Diberikan bilangan bulat positif $k$ dan $n$. Buktikan bahwa

$$
\sum_{i=0}^{k}(-1)^{i}\binom{k}{i}(k-i)^{n}=\begin{cases}n! & \text{jika } k=n, \\ 0 & \text{jika } k>n.\end{cases}
$$


@ID: ON19D2P2U
@BIDANG: Analisis Kompleks
Bilangan bulat positif $n$ dikatakan cantik jika terdapat empat bilangan kompleks $a, b, c,$ dan $d,$ yang memenuhi $a^{n}=b^{n}=c^{n}=d^{n}=1$ dan $a+b+c+d=1$.
(a) Tunjukkan bahwa terdapat bilangan cantik.
(b) Apakah 28 merupakan bilangan cantik? Jelaskan.

@ID: ON19D2P3U
@BIDANG: Struktur Aljabar
Misalkan $R$ suatu gelanggang komutatif dengan unsur kesatuan. Ideal $I$ dan $J$ di $R$ memenuhi $I+J=R$.
(a) Buktikan bahwa pemetaan $\phi:R\rightarrow R/I\times R/J$ yang didefinisikan melalui $\phi(r)=(r+I,r+J)$ merupakan homomorfisma yang surjektif.
(b) Gunakan bagian (a) di atas untuk membuktikan bahwa $R/(I\cap J)$ isomorfik dengan $R/I\times R/J$.

@ID: ON19D2P4U
@BIDANG: Aljabar Linear
Misalkan $A$ matriks tak nol berukuran $n\times(n-k)$, $n<k$ dengan kolom-kolom $A$ adalah himpunan ortogonal yang tidak memuat vektor nol. Misalkan juga $b_{1}, b_{2}, \dots, b_{k}$ adalah vektor kolom berukuran $n-k$. Jika matriks

$$
B=[A_{1}|A_{2}|\dots|A_{n-k}|Ab_{1}|\dots|Ab_{k}]
$$

dengan $A_{i}$ adalah kolom ke-$i$ dari $A,$ tentukan suatu basis bagi ruang nol $B$.

@ID: ON19D2P5U
@BIDANG: Analisis Real
Diberikan barisan bilangan real positif $\{a_{n}\}$ untuk $n\ge0$. Barisan tersebut memenuhi

$$
\sqrt{a_{1}}\ge\sqrt{a_{0}}+1 \quad \text{dan} \quad \left|a_{n+1}-\frac{a_{n}^{2}}{a_{n-1}}\right|\le1,
$$

untuk semua $n$ bilangan bulat positif. Buktikan bahwa barisan $\left\{\frac{a_{n+1}}{a_{n}}\right\}$ konvergen, katakan ke $\alpha$. Selanjutnya buktikan barisan $\left\{\frac{a_{n}}{\alpha^{n}}\right\}$ konvergen.
